% SPDX-FileCopyrightText: 2026 Jonas Whidden
% SPDX-License-Identifier: Apache-2.0 OR CC-BY-4.0
% Generated release copies substitute the two @@...@@ provenance fields.
\documentclass[11pt]{amsart}
\usepackage[T1]{fontenc}
\usepackage{lmodern}
\usepackage[margin=1.05in]{geometry}
\usepackage{amsmath,amssymb,amsthm,booktabs,array}
\usepackage{hyperref}
\hypersetup{colorlinks=true,linkcolor=blue,urlcolor=blue,citecolor=blue,
  pdftitle={A conditional Lean formalization of Erdos Problem 690},
  pdfauthor={Jonas Whidden}}
\newtheorem{theorem}{Theorem}[section]
\newtheorem{lemma}[theorem]{Lemma}
\theoremstyle{definition}
\newtheorem{definition}[theorem]{Definition}
\newcommand{\code}[1]{\path{#1}}
\newcolumntype{L}[1]{>{\raggedright\arraybackslash}p{#1}}
\emergencystretch=2em
\title{A conditional Lean formalization of Erd\H{o}s Problem 690}
\author{Jonas Whidden}
\date{September 2026}
\begin{document}
\begin{abstract}
We give a Lean 4 formalization of the classification of unimodality of the
density of the $k$th smallest distinct prime divisor, conditional on one
explicit estimate for Chebyshev's theta function. The conclusion quantifies
over every positive integer $k$: unimodality holds exactly for $k\leq3$.
The hypothesis is $|\vartheta(x)-x|\leq x/(5\log^2x)$ for all real
$x\geq3{,}594{,}641$, a published estimate of Dusart. Finite prime-count,
short-interval, and theta certificates supply the bounded analytic prefixes;
the finite classification, record-twin certificate, and infinite
Chinese-remainder argument are part of the formal proof, not additional
hypotheses. We describe the proof interfaces, the reduction in analytic
precision needed by the consumer, and the unresolved feasibility of a
complete formal proof of the retained theta estimate. The result is a
conditional formalization, not a claim of an unconditional Lean solution.
\end{abstract}
\maketitle
\begin{center}\small
Copyright \textcopyright\ 2026 Jonas Whidden. Available under the
\href{https://www.apache.org/licenses/LICENSE-2.0}{Apache License 2.0} or
\href{https://creativecommons.org/licenses/by/4.0/}{CC BY 4.0}, at your option.
\end{center}

\section{Problem and contribution}
For a positive integer $n$, order its distinct prime divisors increasingly.
Erd\H{o}s Problem 690 asks about the distribution of a fixed position $k$ in
this list. Cambie studies the associated densities and recurrence
\cite{Cambie}; Wang and Crapis give the complete mathematical answer
\cite{WangCrapis}. We formalize the classification strategy, retaining a
precisely delimited analytic hypothesis from Dusart \cite{Dusart}.
The contribution is the formal construction and composition of the required
proofs and certificates, not a new discovery of the classification.

The distinction between a conditional theorem and an unfinished proof term
is important. Our theorem takes one explicit proposition as an argument.
Every finite classification certificate and every deduction from that
argument is supplied in the theorem's dependency closure. We do not assume
the classification, its finite prefix, or a collection of numerical answers.
An unconditional proof of the analytic argument could be supplied later to
this same function without changing its conclusion.

\section{Definitions and exact statement}
Write $p_0=2<p_1<\cdots$ for the primes, using Lean's zero-based indexing.
For $k\geq1$ and a prime $p$, let
\[
 A_{k,p}=\{n\geq1:p\mid n,\ \#\{q<p:q\text{ prime},\ q\mid n\}=k-1\}.
\]
Its natural density is denoted by $d_k(p)$.
Multiplicity of a prime divisor does not change its position.
\begin{definition}
A sequence $(a_i)_{i\geq0}$ is unimodal if there is an integer $m\geq0$
such that $a_i\leq a_{i+1}$ for $i<m$ and $a_{i+1}\leq a_i$ for $m\leq i$.
Let $\vartheta(x)=\sum_{p\leq x}\log p$, with the sum over primes, and set
$X_0=3{,}594{,}641$. The analytic hypothesis is
\begin{equation}\label{eq:hypothesis}
 \mathcal H:\quad
 \forall x\in\mathbb R,\quad X_0\leq x\ \Longrightarrow\
 |\vartheta(x)-x|\leq\frac{x}{5\log^2x}.
\end{equation}
\end{definition}
\begin{theorem}[Conditional all-$k$ classification]\label{thm:main}
Assume $\mathcal H$. Then for every integer $k\geq1$ the sequence
$(d_k(p_i))_{i\geq0}$ is unimodal if and only if $k\leq3$.
The natural densities defining this sequence exist.
\end{theorem}
The Lean classification declaration is
\begin{quote}\small
\code{PrimeFactorUnimodality.completeClassification_assuming_theta_error}
\end{quote}
with type
\begin{quote}\small
\code{HasThetaLogSquaredError (1 / 5) 3594641}\quad$\to$\quad
\code{CompleteClassification}.
\end{quote}
The latter proposition is exactly
$\forall k\in\mathbb N,\ 1\leq k\Rightarrow
(\operatorname{IsUnimodal}(\operatorname{primeFactorDensity}(k))
\leftrightarrow k\leq3)$.
The separate density-identification theorem is unconditional; it is not an
additional premise of the classification.

\section{Finite densities and adjacent differences}
For a finite prime list $L$, define rational numbers $F(L,r)$ by
\[
\begin{aligned}
 F([],0)&=1,& F([],r+1)&=0,\\
 F(p::L,0)&=\frac{p-1}{p}F(L,0),\\
 F(p::L,r+1)&=\frac{p-1}{p}F(L,r+1)+\frac1pF(L,r).
\end{aligned}
\]
\begin{lemma}[Density formula]
If $L_p$ lists all primes below $p$, then
$d_k(p)=F(L_p,k-1)/p$.
\end{lemma}
\begin{proof}
Divisibility by the primes in $L_p\cup\{p\}$ is periodic modulo their
product. The Chinese remainder theorem counts the residue classes having
exactly $r$ selected divisors. Adding a prime $q$ multiplies the count for
nonselection by $q-1$, while selection contributes the preceding count.
Dividing by the product gives the recurrence. Requiring divisibility by $p$
contributes $1/p$. The density of a periodic set is its residue proportion;
removing the integer zero does not alter the limit.
\end{proof}
This connection is formalized by
\code{hasNaturalDensity_primeFactorDensity}.

For consecutive primes $p<q$ and $r=k-1\geq1$, expansion gives the exact identity
\[
 d_k(q)-d_k(p)
 =\frac{F(L_p,r-1)-(q-p+1)F(L_p,r)}{pq}.
\]
When $F(L_p,r)>0$, the sign is therefore determined by comparing
$F(L_p,r-1)/F(L_p,r)$ with $q-p+1$.
Factoring the nonselection Euler product expresses this ratio as a ratio
of elementary symmetric sums in the positive weights $1/(\ell-1)$ for
primes $\ell<p$. The formal symmetric-sum inequalities bound the ratio in
terms of reciprocal sums and the number of available weights. Positivity
and index restrictions are premises of the corresponding lemmas, rather
than implicit cancellations at zero.

A strict decrease at one index followed by a strict increase at a later
index contradicts unimodality: a mode would have to lie no later than the
first index and strictly after the second. This elementary obstruction is
the common interface between the finite certificates and the uniform tail.

\section{Precise relationship to Wang--Crapis}\label{sec:comparison}
The following comparison refers to version~1 of Wang--Crapis
\cite[Sections~4--6]{WangCrapis}, not to an unspecified later revision.
We retain their density recurrence, symmetric-ratio sign mechanism, and
descent-before-ascent strategy. The arithmetic implementation and the
interface to explicit analytic estimates differ materially.

\begin{center}\small
\begin{tabular}{@{}L{.21\textwidth}L{.34\textwidth}L{.36\textwidth}@{}}
\toprule
Component & Wang--Crapis v1 & This formalization\\
\midrule
Finite/tail split & Finite through $8{,}600{,}001$; tail from
$8{,}600{,}002$ & Finite through $38{,}000$; tail from $38{,}001$\\[4pt]
Large gap & Published gap of length $1{,}113{,}106$ & Only a
$4{,}751$-integer composite sub-block is retained\\[4pt]
Twin ascent & Published $71{,}298$-digit twin & The same twin, with both
primality proofs replayed in Lean\\[4pt]
Later tail interval & Average gap in $(4P,8P]$ & Average gap in
$(4P,9P]$, with corresponding constants reproved\\[4pt]
Reciprocal estimates & Sharp reciprocal error and numerical constants
& Finite Abel lower bound plus an ordinary proved Mertens upper bound\\[4pt]
Analytic trust & Published explicit estimates & One explicit theta
hypothesis; required consequences are derived in Lean\\
\bottomrule
\end{tabular}
\end{center}

In the interval row, $P$ denotes the primorial used by the CRT construction.
The prime-counting interface is also different. Our consumers use
\[
 \pi(n)\geq\frac{n}{\log n}\left(1+\frac1{\log n}\right)
 \quad(n\geq599),\qquad
 \pi(n)\leq\frac{n}{\log n}\left(1+\frac{1.2762}{\log n}\right)
 \quad(n\geq2),
\]
for natural $n$, with real-variable adapters using $\pi(\lfloor x\rfloor)$.
Wang--Crapis Lemma~4.1 instead states the rational-denominator comparisons
$x/(\log x-1)$ and $x/(\log x-1.1)$ on their respective domains
$x>5393$ and $x>60184$. Our altered comparisons are not a claim to prove
those exact inequalities. The short-interval consumer, however, retains
their precise statement: for every real $x\geq3275$, some prime lies in
$(x,x(1+1/(2\log^2x))]$.

The smaller composite block is enough because the new tail begins much
earlier. We need not certify primality of the megagap's endpoints: a proved
composite block already lies inside a gap between actual consecutive
primes. The retained block has one exceptional Fermat check. The retired
full-gap representation required 10,918 such exceptions; its previously
projected 200--250 serial runner-hours are no longer costs of this proof.
This is removal of an obligation by a mathematical replacement, not an
assumption that the larger certificate succeeds.

We also remove two enumeration tasks from the proof closure. A count of at
least 38,000 primes below 500,000 follows from the shared prime-counting
input and $\log(500000)<13.13$, instead of 380 prime-witness modules.
The finite Abel comparison replaces the former 500,001-entry reciprocal
enumeration. Neither change introduces a new hypothesis.

For clarity, the analytic replacement used in the descent is the inequality
\begin{equation}\label{eq:shell}
 S(Y)-S(X)\geq\log\log Y-\log\log X
       -\frac{6381}{5000\log^2X}-\frac1X,
 \qquad 600\leq X\leq Y,
\end{equation}
where $X,Y$ are integers and $S(t)=\sum_{p\leq t}1/p$.
The $1/X$ term is retained from the floor correction in the integral
prime-counting lower bound. In the CRT descent we use
$r=k-1$, $X=p_{r-1}-1$ (in zero-based indexing) and
$Y=P=\prod_{p\leq q}p$, where the selected prime $q$ satisfies
\[
 \frac{491r}{500\log r}<q<
 \frac{2519}{2500}\frac{491r}{500\log r}.
\]
A tangent bound gives a strictly positive affine lower margin for all
$r\geq38{,}000$, with rational intercept
$877288074251/1231383407380000$ and slope $3078619/10355190$.
Thus the lowered cutoff is proved uniformly, not inferred by sampling.

For the ascent we use the established ordinary Mertens estimate and the
coarse bound on its constant $B<2/3$. We do not need the published decimal
enclosures of $B$ and $\sum_p1/(p(p-1))$ used in the original argument.
The declarations in \code{PrimeCountingReciprocalShell},
\code{TailShellNumerics} and \code{UniformTailClosedMertens}
make these replacements explicit. Our all-$k$ conclusion is unchanged;
the paper's unconditional mathematical theorem and our conditional formal
theorem have different trust boundaries.

\section{Finite and infinite classification ranges}
The finite classification runs through $k=38{,}000$.
The first three cases follow from exact recurrence inequalities.
For $4\leq k\leq38{,}000$, the development combines initial exact valleys,
symmetric-ratio estimates, a certified record-gap descent and a later
certified twin-prime ascent. The record twins are actual primality
certificates in the proof closure. They are not asserted as trusted data.
The assembled finite theorem takes prime-counting bounds as its sole
analytic input, supplied below by $\mathcal H$.

An Abel-summation estimate removes the need to enumerate a long reciprocal
prime list in the upper record range. The formal comparison bounds the
endpoint loss by $1/50$ for $k\geq10{,}372$ and obtains a reciprocal-sum
margin greater than eight. The gap multiplier then gives
$4753\cdot8=38{,}024>k-1$ throughout the remaining finite range.
This replacement removes the old large reciprocal enumeration from the
conditional proof's import closure.

For every $k\geq38{,}001$, a Chinese-remainder construction supplies a
composite block and hence a descent. Bounds on the primorial control the
size of the construction. Prime-counting and short-interval estimates
provide a subsequent consecutive-prime pair, for which the symmetric-ratio
comparison gives an ascent. This is a theorem for an arbitrary $k$ in the
tail, not a finite search. The finite and infinite cases meet without a gap.

\section{Deriving the analytic interfaces from one hypothesis}
Three proved adapters consume $\mathcal H$:
\begin{enumerate}
\item \code{hasDusartPrimeCountingBounds_of_logSquared_ray};
\item \code{hasDusartShortIntervalPrime_of_logSquared_ray};
\item \code{hasTailThetaBounds_of_logSquared_ray}.
\end{enumerate}
The first uses prime-log Abel summation, differential comparisons and an
exact anchor count $\pi(X_0)=256{,}357$. Bounds for $\log X_0$ and the same
hypothesis $\mathcal H$ prove the two required anchor comparisons. The
second derives the large-input short-interval conclusion by comparing
theta at the interval endpoints. Finite prefixes supply the missing small
inputs in both adapters.

The third adapter uses only the precision required by the tail construction:
for natural $q\geq3501$,
\begin{equation}\label{eq:tail}
 q\left(1-\frac{1.2323}{\log q}\right)
 <\vartheta(q)<1.001q.
\end{equation}
For $q\geq X_0$, $\log q>15$ implies
$1/(5\log^2q)<1/1000$ and
$1/(5\log^2q)<1.2323/\log q$; equation~\eqref{eq:hypothesis}
then proves~\eqref{eq:tail}. The interval $3501\leq q<X_0$ is certified
separately. In the primorial comparison,
$1.0076\cdot1.001+0.0007=1.0093076<1.01$ preserves the required margin.
Consequently the stronger published global upper coefficient
$1+1/36{,}260$ need not be independently formalized for this classification.
We do not assert that the weaker interface proves that stronger bound.

\begin{proof}[Assembly of Theorem~\ref{thm:main}]
Apply the first adapter to $\mathcal H$ and supply its result to the finite
classification theorem
\begin{quote}\small
\code{completeClassification_through38000_of_primeCounting}.
\end{quote}
This discharges the finite classification, including its record-twin input.
The second and third adapters give the remaining hypotheses of the uniform
tail theorem. The final reduction is
\begin{quote}\small
\code{completeClassification_of_finite_prefix_and_tailTheta}.
\end{quote}
It splits at $38{,}000$. The finite branch gives the equivalence; the infinite
branch excludes unimodality and also has $k>3$. Density existence follows
from the periodic counting lemma. No further analytic parameter remains.
\end{proof}

\section{Finite certificate representation}
Numerical search and row generation may happen outside Lean, but their
outputs are consumed by proved checking and coverage theorems. Compact rows
share assemblers rather than creating a theorem module for each prime.
The prime-count prefix above 1000 contains 2,093 rows in nine batches;
the short-interval prefix from 3275 contains 2,022 rows in eight batches.
The theta prefix uses 9,765 transitions in 26 batches. All end at the
published analytic cutoff $X_0$ with the required endpoint conventions.

Coverage is part of the proof. A prime-count row ending at a prime uses the
count just before the jump for its half-open cell; the next cell includes
the jump. A short-interval witness must be strictly larger than its input,
including at shared boundaries. Theta rows combine exact count increments
with rational logarithm bounds, and the assembly checks adjacency, terminal
coverage and strict margins. A table with omitted cells is not a proof.
The exact anchor count uses a certified packed sieve. External primality
tests or floating-point estimates do not discharge a Lean hypothesis.

\section{What remains in the Dusart formalization}
Equation~\eqref{eq:hypothesis} is a published mathematical result, not a
conjecture introduced here. Its complete formal derivation remains outside
this theorem. The obstacle is reproducing its analytic and finite numerical
ingredients within a justified computational budget, not evidence against
Dusart's theorem or against the possibility of formalizing it.

Several reductions improve the situation. The weaker consumer-specific
theta interface avoids a separate stronger-bound requirement extending
into the trillions. The finite analytic prefixes and Abel anchors leave no
extra assumptions. Packed sieves and compact count/log rows reduce the
cost of finite prime information. Reflected explicit-formula estimates and
finite-height zero information offer candidate routes for the remaining
unbounded estimate, but their complete numerical obligations still matter.

In particular, approximate zero locations do not certify all zeros below a
height. A rigorous approach needs validated function enclosures, signs or
other root certificates, disjoint ordered brackets, and a completeness
argument such as a verified zero-count comparison. Gamma factors,
truncation errors, rounding bounds and coverage must be included. An
external fast sign evaluation does not itself produce these formal facts.
Shared transforms and convolution-based evaluators are promising, but their
input authenticity and complete error budgets cannot be replaced by timing
one integer multiplication or one endpoint comparison.

\subsection{Measured costs and rejected approaches}
The development host has two AArch64 Neoverse-N1 cores. The figures below
are single-run component wall times, not confidence intervals, CPU billing
or predictions for another runner. Warm dependency imports are included
where a complete Lean process is timed; external generation is excluded
unless explicitly stated. Sampled resident-memory readings can miss peaks.
The companion \code{paper/conditional-resources.json} records the selected
pilot results, their scope and hashes of the original measurement reports.
Missing experiment-revision information is identified rather than inferred.

\begin{center}\small
\begin{tabular}{@{}L{.35\textwidth}L{.16\textwidth}L{.40\textwidth}@{}}
\toprule
Work measured & Wall time & Interpretation\\
\midrule
Packed count $\pi(X_0)$ & $5.09$ s & Exact count and saved Lean artifact\\
Short-interval prefix & $75.7$ s & Eight batches combined\\
Weak theta prefix & $700$ s & Twenty-six batches combined\\
Packed count $\pi(30{,}000{,}000)$ & $27.52$ s & Count only, not a theta bound\\
Low-height Euler sign/count block & $128.25$ s & Complete local block, not
global zero completeness\\
Packed real convolution, length $65{,}536$ & $32.81$ s & Nonnegative
96-bit coefficients; no saved artifact or complex error semantics\\
\bottomrule
\end{tabular}
\end{center}

Individual primality certificates did not solve the bulk counting problem.
For example, 64 Pocklington targets near $10^{10}$ took 6.81 seconds,
including process/import cost; shared Pratt certificates took 8.97 seconds.
The problem is not that one prime is intrinsically expensive. One needs
complete counting or weighted-sum coverage over a large domain, with all
reuse represented efficiently in the proof. A BPSW probable-prime test alone
does not certify primality, and primality certificates alone do not prove
that a list contains every relevant prime. ECPP was considered as a
certificate alternative; no implemented end-to-end ECPP speedup is claimed.

A Meissel-style census for a candidate count mesh through $3\cdot10^9$
modeled 307,016,043 rows. At a measured 422.33 checked rows per second,
checking contributed 8.41 serial days; estimated row-data elaboration added
2.07 days, with 13.87 GiB of source. Data elaboration here is Lean's cost of
loading the generated mathematical data, not the external script's run
time. The roughly 252-worker-hour total omitted shared-table proofs,
logarithmic bounds, theta reconstruction and zero verification. It was a
finite counting scenario, not a previously complete ten-day Dusart proof.
Packed sieve counts improve this component, but neither one count nor a
linear extrapolation of the 30-million pilot prices the full three-billion
theta band.

Direct high-order Euler evaluation was also screened. Shared logarithms
help, but moments still change between windows; a design with a long
Euler prefix and a full radix trace creates enormous arithmetic and data
volumes at high height. The measured low-height sign/count block therefore
does not justify a bulk rollout. We rejected that layout for high-height
verification, rather than extrapolating its first successful block as a
feasible full proof.

\subsection{Analytic height reduction and ongoing zero work}
The relevant zero height is not immutable. Stronger reflected zero-sum
envelopes and smoothing comparisons have reduced the screened height from
the initial $10^9$ scale to a candidate $H=170{,}000{,}000$. This candidate
still requires its finite analytic caps, continuous coverage and low-zero
completeness hypotheses. Smaller-height experiments did not supply a
replacement proof: a tested Stieltjes-envelope design at height
$25{,}000{,}000$ attained a normalized error near $0.855$ in its difficult
region, above the required $0.2$. Failure of that parameter search does
not rule out all lower-height analytic methods.

At $H=170{,}000{,}000$, the zero-count main term suggests roughly
436 million positive zeros; this is a workload estimate, not a certified
count. External 96-bit Arb evaluations of 17 nearby endpoint signs took
about 0.0383 seconds, but omitted Lean checking and completeness. Even
fast numerical signs and a supplied zero list do not eliminate the cost of
authenticating those values and proving that no zeros were missed.

A screened Gaussian-window/convolution scenario uses 41,504 windows and
15 Taylor columns, hence 622,560 complex convolutions. Its dense layout
has about $1.632\cdot10^{11}$ complex input cells: at two 96-bit components
per cell, about 3.92 TB before outputs, syntax and metadata. That layout
was rejected. Merely charging each complex convolution as one smaller
real pilot already gives $622560\cdot32.81/3600\approx5674$ worker-hours.
This deliberately incomplete proxy is not a valid total runtime forecast
or a lower bound for all convolution algorithms; it demonstrates why the
particular pilot cannot support a days-scale claim for the whole system.

Ongoing formal infrastructure includes arbitrary-point Euler expansions,
Gamma asymptotics with rational input interfaces, phase-projection sign
criteria, rational radix-transform checking, and shared rounded-product
and correlation error assembly. These prove semantic reductions and error
propagation. They do not yet provide all concrete elementary-function
enclosures, high-height sign margins, exceptional cases or the final
zero-count comparison. Shared transforms and compact reconstruction may
avoid the rejected dense representation; they require a measured complete
sign-producing block and a full-domain cost census before selection.

No measured feasible complete route for discharging $\mathcal H$ is claimed
here. Parallel execution can reduce elapsed time only after independence,
per-worker memory and total work have been established. Checkpointing
preserves completed work but does not make an infeasible computation
feasible. A future improvement must account for the whole path, including
generation, formal checking, storage, import costs and completeness.
For $N$ independent blocks costing $t$ worker-seconds and $m$ bytes each,
$W$ memory-admissible workers have an ideal compute floor $Nt/W$, before
communication, startup and dependencies; admissibility requires the
aggregate working sets to fit available memory. Our operational budget
permits days overall with jobs of at most six hours. Serial dependency
chains prevent simultaneous large imports, and content-addressed
checkpoints preserve completed units across unrelated commits. These are
resource controls, not additional mathematical assumptions.

\section{Trust boundary and reproducibility}
The conditional theorem uses the usual Lean axioms
\code{propext}, \code{Classical.choice} and \code{Quot.sound}.
The retained analytic premise is an explicit theorem argument, not a
project-local axiom. The certificate closure uses kernel-checked proofs;
native evaluation is not used to establish this theorem. The assertion
proved is $\mathcal H\Rightarrow\mathrm{CompleteClassification}$, not
$\mathrm{CompleteClassification}$ alone.

The version accompanying this paper identifies the exact theorem type,
axiom report, dependency locks and source-addressed certificate units.
Compiled artifacts are conveniences for reuse, not replacements for the
source or the explicit hypothesis. The unconditional Challenge and Solution
are separate entry points and are not certified by this conditional result.
The original three-bound conditional entry point remains available for
comparison, but this paper concerns only the single-theta theorem.

\begin{center}\small
Source commit: \texttt{@@COMMIT@@}\\
Lean toolchain: \texttt{@@TOOLCHAIN@@}
\end{center}

\section{Conclusion}
The classification for every positive $k$ is reduced in Lean to the single
explicit estimate~\eqref{eq:hypothesis}, with its finite and infinite proof
components retained. This isolates a reusable, fully specified conditional
result while leaving the mathematical and computational burden of the
unbounded theta estimate visible. Removing that argument requires its
formal proof, not a weaker classification statement or an expanded hidden
trust boundary.

\begin{thebibliography}{9}
\bibitem{Cambie} Stijn Cambie, \emph{Resolution of Erd\H{o}s' problems about
unimodularity}, arXiv:2501.10333, 2025.
\url{https://arxiv.org/abs/2501.10333}.
\bibitem{WangCrapis} Shouqiao Wang and Davide Crapis,
\emph{A Complete Answer to Erd\H{o}s Problem 690}, arXiv:2605.08542v1, 2026.
\url{https://arxiv.org/abs/2605.08542v1}.
\bibitem{Dusart} Pierre Dusart, \emph{Estimates of Some Functions over Primes
without R.H.}, arXiv:1002.0442v1, 2010.
\url{https://arxiv.org/abs/1002.0442v1}.
\end{thebibliography}
\end{document}
